\subsection{里斯自同态的扰动}

设 E 是巴拿赫空间。回忆（参见 III, p. 5, 命题 3）：E 的紧自同态所成的集 \(\mathcal{L}^c (\mathrm{E})\) 构成巴拿赫代数 \(\mathcal{L}(\mathrm{E})\) 的一个闭双边理想。商巴拿赫代数称为 E 的卡尔金代数；记作 \(\mathcal{Calk}(\mathrm{E})\) 。我们用 \(\pi\) 表示从 \(\mathcal{L}(\mathrm{E})\) 到 \(\mathcal{Calk}(\mathrm{E})\) 上的典范同态。由 III, p. 74 的推论 1，E 的自同态 \(u\) 是弗雷德霍姆自同态，当且仅当 \(\pi (u)\) 在代数 \(\mathcal{Calk}(\mathrm{E})\) 中可逆。

命题 2. --- 设 \(u \in \mathcal{L}(\mathrm{E})\) 满足 \(\|1_{\mathrm{E}} - \pi(u)\| < 1\) （在代数 \(\mathcal{Calk}(\mathrm{E})\) 中）。则 \(u\) 是 E 的里斯自同态。

设 \(r \geqslant 1\) 是使 \(\| 1_{\mathrm{E}} - \pi(u)\|^{r} < \frac{1}{2}\) 的整数。设 \(\mathrm{P} \in \mathbf{C}[\mathrm{X}]\) 是多项式 \(\frac{1 - (1 - \mathrm{X})^r}{\mathrm{X}}\) 。令 \(v = 1_{\mathrm{E}} - (1_{\mathrm{E}} - u)^r\) 。于是我们有 \(v = u\mathrm{P}(u)\) 与 \(\| 1_{\mathrm{E}} - \pi(v)\| < \frac{1}{2}\) 。由于 E 的自同态 \(u\) 与 \(\mathrm{P}(u)\) 交换，只需证明 \(v\) 是 E 的里斯自同态（III, p. 49, 命题 9）。

由于 \(\|1_{\mathrm{E}}-\pi(v)\|<1/2\) ，由空间 \(\mathcal{Calk}(\mathrm{E})\) 中商范数的定义，存在 E 的一个紧自同态 h 和 E 的一个自同态 w，使得 \(v=1_{E}+h+w\) 且 \(\|w\|<\frac{1}{2}\) 。由 I, p. 22 的推论 1，元素 \(1_{E}+w\) 是 E 的一个自同构。由于 h 是紧的，\(v=(1_{\mathrm{E}}+w)+h\) 是 E 的指标为 0 的弗雷德霍姆自同态（III, p. 74 的推论 2）。对每个整数 \(n\geqslant0\) ，用 \(N_{n}\) 表示 \(v^{n}\) 的核。为证明 v 是 E 的里斯自同态，只需证明存在一个整数 \(n \geqslant 0\) 使 \(N_{n} = N_{n+1}\) （III, p. 46, 定义 2 及 III, p. 46, 注记）。

我们用反证法，假设序列 \((\mathrm{N}_{n})\) 严格递增。对每个 \(n \in N\) ，用 \(p_{n}\) 表示从 E 到赋范空间 \(E/N_{n}\) 上的典范映射。设 c 是使 \(2\|w\| < c < 1\) 的实数。设 \(n \in N\) 。由于 \(N_{n+1}\) 异于 \(N_{n}\) ，存在一个元素 \(x_{n} \in N_{n+1}\) ，使得 \(\|p_{n}(x_{n})\| = c\) 且 \(\|x_{n}\| < 1\) （事实上，存在 \(y \in N_{n+1}/N_{n}\) 范数为 c，故对每个 \(\varepsilon > 0\) ，存在 \(x_{n} \in N_{n+1}\) 使得 \(p_{n}(x_{n}) = y\) 且 \(\|x_{n}\| \leqslant c + \varepsilon\) ）。

设 \(m\) 、\(n\) 是使 \(m > n \geqslant 0\) 的两个整数。我们有

\[
\begin{array}{r l} & {\| h (x _ {m}) - h (x _ {n}) \| = \| v (x _ {m} - x _ {n}) - (1 _ {\mathrm{E}} + w) (x _ {m} - x _ {n}) \|} \\ & {\qquad \geqslant \| (v - 1 _ {\mathrm{E}}) (x _ {m} - x _ {n}) \| - \| w (x _ {m} - x _ {n}) \|} \\ & {\qquad \geqslant \| (v - 1 _ {\mathrm{E}}) (x _ {m} - x _ {n}) \| - 2 \| w \|} \end{array}\tag{1}
\]

因为 \(\| x_m\| \leqslant 1\) 且 \(\| x_n\| \leqslant 1\) 。此外，注意到

\[
(v - 1 _ {\mathrm{E}}) (x _ {m} - x _ {n}) = v (x _ {m} - x _ {n}) + x _ {n} - x _ {m}.
\]

但由于 n \textless{} m 且 \(v(\mathrm{N}_{m+1}) \subset \mathrm{N}_{m}\) ，我们有 \(v(x_{m} - x_{n}) + x_{n} \in \mathrm{N}_{m}\) ，由此

\[
\left\| v (x _ {m} - x _ {n}) + x _ {n} - x _ {m} \right\| \geqslant \left\| p _ {m} (x _ {m}) \right\| = c.
\]

于是不等式 (1) 给出下界

\[
\left\| h \left(x _ {m}\right) - h \left(x _ {n}\right) \right\| \geqslant c - 2 \| w \| > 0.\tag{2}
\]

于是序列 \((h(x_{m}))_{m\in\mathbf{N}}\) 在 E 中没有聚点；由于 h 是紧的（III, p. 2 的注记 1），E 的单位球在 h 下的像是相对紧的，这就产生矛盾。

推论 1. --- 设 E 是分离局部凸空间，且 \(u \in \mathcal{L}(\mathrm{E})\) 。为使 u 是里斯自同态，必须且只需存在 \(\mathcal{L}(\mathrm{E})\) 的一个元素 v，它与 u 交换且使 \(1_{\mathrm{E}} - u \circ v\) 是紧的。

由于每个有限秩连续线性映射都是紧的，该条件是必要的（III, p. 47, 命题 6 d)）。我们来证明它是充分的。设 v 是 \(\mathcal{L}(\mathrm{E})\) 的一个元素，它与 u 交换，且使 E 的自同态 \(h = 1_{E} - u \circ v\) 是紧的。存在一个巴拿赫空间 G、一个连续线性映射 \(p: E \to G\) 和一个紧线性映射 \(q: G \to E\)，使得 \(h = q \circ p\) （III, p. 5, 命题 5）。G 的自同态 \(p \circ q\) 是紧的。由命题 2，线性映射 \(1_{G} - p \circ q\) 是 G 的里斯自同态。但此时 \(u \circ v = 1_{E} - q \circ p\) 是 E 的里斯自同态（III, p. 49, 命题 10, g)）。由于 u 与 v 交换，u 是 E 的里斯自同态（III, p. 49, 命题 9）。

推论 2. --- 设 E 是分离局部凸空间，u 是 E 的里斯自同态，\(h \in \mathcal{L}(\mathrm{E})\) 是紧线性映射。若 u 与 h 交换，则 \(u + h\) 是 E 的里斯自同态。

设 \(v\) 是 \(u\) 的典范拟逆。它与 \(u\) 交换，也与 \(h\) 交换（III, p. 47, 命题 6），从而与 \(u + h\) 交换。E 的自同态 \(1_{\mathrm{E}} - (u + h) \circ v\) 是有限秩连续线性映射 \(1_{\mathrm{E}} - u \circ v\) 与紧线性映射 \(-h \circ v\) 之和，故是紧的。由此得出 \(u + h\) 是 E 的里斯自同态（推论 1）。

命题 3. --- 设 E 是巴拿赫空间。设 u 是 E 的一个自同态，A 是它在 \(\mathcal{L}(\mathrm{E})\) 中的交换子。\(\pi(\mathrm{A})\) 在代数 \(\mathcal{Calk}(\mathrm{E})\) 中的闭包 B 是一个巴拿赫代数。为使 u 是 E 的里斯自同态，必须且只需 \(\pi(u)\) 在 B 中可逆。

由于 \(\pi(A)\) 是 \(\mathcal{Calk}(E)\) 的一个含单位元的子代数，其闭包 B 是一个巴拿赫代数。若 \(u\) 是 E 的里斯自同态，则元素 \(\pi(u)\) 在 \(\pi(A)\) 中有一个逆（推论 1），从而在 B 中有逆。反之，假设 \(\pi(u)\) 在 B 中有一个逆 \(s\) 。由 B 的定义，存在 A 的一个元素 \(v\) 使得

\[
\left\| s - \pi (v) \right\| <   \frac {1}{\left\| \pi (u) \right\|},
\]

由此 \(\|1_{E}-\pi(u\circ v)\|<1\) 。由命题 2，线性映射 \(u\circ v\) 是 E 的里斯自同态。由于 u 与 v 交换，u 也是（III, p. 49, 命题 9）。

